\documentclass[aspectratio=169,11pt]{beamer}

\usetheme{Madrid}
\usecolortheme{wolverine}
\IfFileExists{presentation/theme/mihajlo-elfak-logos.sty}{%
  \NeedsTeXFormat{LaTeX2e}
\ProvidesFile{mihajlo-elfak-logos.sty}[2026/06/12 ELFak corner logos for Beamer]

\RequirePackage{graphicx}
\RequirePackage{tikz}

\IfFileExists{presentation/theme/mihajlo-elfak-university-logo.png}{%
  \newcommand{\mihajloelfakassetpath}{presentation/theme/}%
}{%
  \newcommand{\mihajloelfakassetpath}{theme/}%
}

\definecolor{ElfakPrimary}{HTML}{164D7D}
\definecolor{ElfakPrimaryDark}{HTML}{0E3252}
\definecolor{ElfakPrimaryLight}{HTML}{2F75B5}
\definecolor{ElfakAccent}{HTML}{A51003}
\definecolor{ElfakAccentDark}{HTML}{730B02}
\definecolor{ElfakLight}{HTML}{E4EDF5}

\setbeamercolor{structure}{fg=ElfakPrimary}
\setbeamercolor{title}{fg=white,bg=ElfakPrimary}
\setbeamercolor{subtitle}{fg=white}
\setbeamercolor{frametitle}{fg=white,bg=ElfakPrimary}
\setbeamercolor{palette primary}{fg=white,bg=ElfakPrimary}
\setbeamercolor{palette secondary}{fg=white,bg=ElfakPrimaryDark}
\setbeamercolor{palette tertiary}{fg=white,bg=ElfakAccent}
\setbeamercolor{palette quaternary}{fg=white,bg=ElfakAccentDark}
\setbeamercolor{block title}{fg=white,bg=ElfakPrimary}
\setbeamercolor{block body}{fg=black,bg=ElfakLight}
\setbeamercolor{alerted text}{fg=ElfakAccent}
\setbeamercolor{itemize item}{fg=ElfakAccent}
\setbeamercolor{itemize subitem}{fg=ElfakPrimary}
\setbeamercolor{enumerate item}{fg=ElfakAccent}
\setbeamercolor{author in head/foot}{fg=white,bg=ElfakPrimaryDark}
\setbeamercolor{title in head/foot}{fg=white,bg=ElfakPrimary}
\setbeamercolor{date in head/foot}{fg=white,bg=ElfakPrimaryLight}

\newcommand{\mihajloelfaklogos}{%
  \begin{tikzpicture}[remember picture,overlay]
    \node[anchor=north west,inner sep=0pt,xshift=0.08cm,yshift=-0.05cm]
      at (current page.north west)
      {\includegraphics[width=0.95cm]{\mihajloelfakassetpath mihajlo-elfak-university-logo.png}};
    \node[anchor=north east,inner sep=0pt,xshift=-0.10cm,yshift=-0.06cm]
      at (current page.north east)
      {\includegraphics[width=0.82cm]{\mihajloelfakassetpath mihajlo-elfak-faculty-logo.png}};
  \end{tikzpicture}%
}

\AddToHook{shipout/foreground}{%
  \ifnum\insertframenumber=1
    \mihajloelfaklogos
  \fi
}

\setbeamertemplate{footline}{%
  \ifnum\insertframenumber>1
    \leavevmode%
    \hbox{%
      \begin{beamercolorbox}[wd=.28\paperwidth,ht=2.5ex,dp=1.1ex,leftskip=.8em]{author in head/foot}%
        \usebeamerfont{author in head/foot}%
        \usebeamercolor[fg]{author in head/foot}%
        \NoHyper\insertshortauthor\endNoHyper
      \end{beamercolorbox}%
      \begin{beamercolorbox}[wd=.44\paperwidth,ht=2.5ex,dp=1.1ex,center]{title in head/foot}%
        \usebeamerfont{title in head/foot}%
        \usebeamercolor[fg]{title in head/foot}%
        \NoHyper\insertshortsubtitle\endNoHyper
      \end{beamercolorbox}%
      \begin{beamercolorbox}[wd=.28\paperwidth,ht=2.5ex,dp=1.1ex,rightskip=.8em plus1fil]{date in head/foot}%
        \usebeamerfont{date in head/foot}%
        \usebeamercolor[fg]{date in head/foot}%
        \NoHyper\insertshortinstitute\endNoHyper\hspace{1em}%
        \insertframenumber/\inserttotalframenumber
      \end{beamercolorbox}%
    }%
  \fi
}
%
}{%
  \NeedsTeXFormat{LaTeX2e}
\ProvidesFile{mihajlo-elfak-logos.sty}[2026/06/12 ELFak corner logos for Beamer]

\RequirePackage{graphicx}
\RequirePackage{tikz}

\IfFileExists{presentation/theme/mihajlo-elfak-university-logo.png}{%
  \newcommand{\mihajloelfakassetpath}{presentation/theme/}%
}{%
  \newcommand{\mihajloelfakassetpath}{theme/}%
}

\definecolor{ElfakPrimary}{HTML}{164D7D}
\definecolor{ElfakPrimaryDark}{HTML}{0E3252}
\definecolor{ElfakPrimaryLight}{HTML}{2F75B5}
\definecolor{ElfakAccent}{HTML}{A51003}
\definecolor{ElfakAccentDark}{HTML}{730B02}
\definecolor{ElfakLight}{HTML}{E4EDF5}

\setbeamercolor{structure}{fg=ElfakPrimary}
\setbeamercolor{title}{fg=white,bg=ElfakPrimary}
\setbeamercolor{subtitle}{fg=white}
\setbeamercolor{frametitle}{fg=white,bg=ElfakPrimary}
\setbeamercolor{palette primary}{fg=white,bg=ElfakPrimary}
\setbeamercolor{palette secondary}{fg=white,bg=ElfakPrimaryDark}
\setbeamercolor{palette tertiary}{fg=white,bg=ElfakAccent}
\setbeamercolor{palette quaternary}{fg=white,bg=ElfakAccentDark}
\setbeamercolor{block title}{fg=white,bg=ElfakPrimary}
\setbeamercolor{block body}{fg=black,bg=ElfakLight}
\setbeamercolor{alerted text}{fg=ElfakAccent}
\setbeamercolor{itemize item}{fg=ElfakAccent}
\setbeamercolor{itemize subitem}{fg=ElfakPrimary}
\setbeamercolor{enumerate item}{fg=ElfakAccent}
\setbeamercolor{author in head/foot}{fg=white,bg=ElfakPrimaryDark}
\setbeamercolor{title in head/foot}{fg=white,bg=ElfakPrimary}
\setbeamercolor{date in head/foot}{fg=white,bg=ElfakPrimaryLight}

\newcommand{\mihajloelfaklogos}{%
  \begin{tikzpicture}[remember picture,overlay]
    \node[anchor=north west,inner sep=0pt,xshift=0.08cm,yshift=-0.05cm]
      at (current page.north west)
      {\includegraphics[width=0.95cm]{\mihajloelfakassetpath mihajlo-elfak-university-logo.png}};
    \node[anchor=north east,inner sep=0pt,xshift=-0.10cm,yshift=-0.06cm]
      at (current page.north east)
      {\includegraphics[width=0.82cm]{\mihajloelfakassetpath mihajlo-elfak-faculty-logo.png}};
  \end{tikzpicture}%
}

\AddToHook{shipout/foreground}{%
  \ifnum\insertframenumber=1
    \mihajloelfaklogos
  \fi
}

\setbeamertemplate{footline}{%
  \ifnum\insertframenumber>1
    \leavevmode%
    \hbox{%
      \begin{beamercolorbox}[wd=.28\paperwidth,ht=2.5ex,dp=1.1ex,leftskip=.8em]{author in head/foot}%
        \usebeamerfont{author in head/foot}%
        \usebeamercolor[fg]{author in head/foot}%
        \NoHyper\insertshortauthor\endNoHyper
      \end{beamercolorbox}%
      \begin{beamercolorbox}[wd=.44\paperwidth,ht=2.5ex,dp=1.1ex,center]{title in head/foot}%
        \usebeamerfont{title in head/foot}%
        \usebeamercolor[fg]{title in head/foot}%
        \NoHyper\insertshortsubtitle\endNoHyper
      \end{beamercolorbox}%
      \begin{beamercolorbox}[wd=.28\paperwidth,ht=2.5ex,dp=1.1ex,rightskip=.8em plus1fil]{date in head/foot}%
        \usebeamerfont{date in head/foot}%
        \usebeamercolor[fg]{date in head/foot}%
        \NoHyper\insertshortinstitute\endNoHyper\hspace{1em}%
        \insertframenumber/\inserttotalframenumber
      \end{beamercolorbox}%
    }%
  \fi
}
%
}
\usepackage{fontspec}
\setsansfont{Latin Modern Sans}
\usepackage[provide=*,serbian-latin]{babel}

\title{Projektovanje i implementacija AI MVP}
\subtitle[ASAP]{Automatska Semantička Analiza Proizvoda}
\author[Mihajlo Madić 2119]{Mihajlo Madić 2119}
\institute[Inteligentni sistemi]{Inteligentni sistemi}
\date{2026}

\newcommand{\planned}{\textcolor{gray}{\footnotesize Planirano --- dopuniti tokom razvoja}}

\begin{document}

\begin{frame}
  \titlepage
\end{frame}

\begin{frame}{Problem i ideja}
  \begin{itemize}
    \item Podaci o proizvodima su rasuti, a korisniku su potrebni odmah, u trenutku kupovine.
    \item ASAP povezuje skenirani barkod sa metapodacima proizvoda.
    \item Semantičke reprezentacije omogućavaju pretragu sličnih proizvoda i personalizovane preporuke.
  \end{itemize}
\end{frame}

\begin{frame}{Planirane AI funkcionalnosti}
  \begin{enumerate}
    \item lokalno prepoznavanje i dekodiranje barkoda;
    \item embedding reprezentacije proizvoda;
    \item top-N semantička pretraga kosinusnom sličnošću;
    \item opcioni MMR za raznovrsnije rezultate;
    \item korisnički profil kao centroid prethodnih interakcija;
    \item opciona klasterizacija i PCA vizuelna analitika.
  \end{enumerate}
\end{frame}

\begin{frame}{Početni tok sistema}
  \centering
  \begin{columns}[T,onlytextwidth]
    \column{0.30\textwidth}
      \begin{block}{Mobilna aplikacija}
        Kamera $\rightarrow$ barkod\\[0.4em]
        Prikaz proizvoda i preporuka
      \end{block}
    \column{0.05\textwidth}
      \centering\vspace{1.4cm}$\Longleftrightarrow$
    \column{0.30\textwidth}
      \begin{block}{Backend}
        API servis\\[0.4em]
        Semantička pretraga i preporuke
      \end{block}
    \column{0.05\textwidth}
      \centering\vspace{1.4cm}$\Longleftrightarrow$
    \column{0.30\textwidth}
      \begin{block}{Podaci}
        Metapodaci proizvoda\\[0.4em]
        Vektorski indeks
      \end{block}
  \end{columns}
\end{frame}

\section{Definisanje problema}

\begin{frame}{Ciljni korisnici i konkurentska prednost}
  \planned
  \begin{itemize}
    \item Ko su primarni korisnici?
    \item Koji problem rešavaju danas i kojim alternativama?
    \item Po čemu je ASAP merljivo bolji ili jednostavniji?
  \end{itemize}
\end{frame}

\section{Tehnologije i podaci}

\begin{frame}{Izabrane tehnologije}
  \planned
  \begin{itemize}
    \item mobilna platforma i programski jezik;
    \item backend i skladište podataka;
    \item ML Kit i model za embedding reprezentacije;
    \item vektorski indeks i izvori podataka o proizvodima.
  \end{itemize}
\end{frame}

\section{Implementacija}

\begin{frame}{MVP demonstracija}
  \planned
  \begin{enumerate}
    \item skeniranje barkoda;
    \item pribavljanje podataka o proizvodu;
    \item prikaz semantički sličnih proizvoda;
    \item personalizovana preporuka.
  \end{enumerate}
\end{frame}

\begin{frame}{Odluke, problemi i korekcije}
  \planned
  \begin{itemize}
    \item ključne tehničke odluke i kompromisi;
    \item povratne informacije sa MVP prezentacije;
    \item izmene sprovedene nakon prezentacije.
  \end{itemize}
\end{frame}

\section{Evaluacija}

\begin{frame}{Rezultati testiranja}
  \planned
  \begin{itemize}
    \item uspešnost i vreme skeniranja;
    \item kvalitet semantičke pretrage i preporuka;
    \item povratne informacije potencijalnih korisnika;
    \item ograničenja i sledeći koraci.
  \end{itemize}
\end{frame}

\begin{frame}{Zaključak}
  \planned
  \begin{itemize}
    \item ostvarene funkcionalnosti u odnosu na plan;
    \item najvažniji rezultat projekta;
    \item mogući pravci daljeg razvoja.
  \end{itemize}
\end{frame}

\end{document}
